% Part: first-order-logic
% Chapter: completeness
% Section: maximally-consistent-sets

% Definition of maximally consistent sets. Properties of mcs required
% for completeness proved are in maximally-consistent-sets.tex in the
% chapter on the proof system used.

\documentclass[../../../include/open-logic-section]{subfiles}

\begin{document}

\olfileid{fol}{com}{mcs}
\olsection{\usetoken{P}{sentence}ના મહત્તમ સુસંગત ગણો}

\begin{defn}[મહત્તમ સુસંગત ગણ]
\ollabel{mcs}
!!{sentence}sનો ગણ~$\Gamma$ \emph{મહત્તમ સુસંગત} છે જો અને માત્ર જો
\begin{enumerate}
\item $\Gamma$ સુસંગત હોય, અને
\item જો $\Gamma \subsetneq \Gamma'$, તો $\Gamma'$ અસુસંગત હોય.
\end{enumerate}
\end{defn}

\begin{explain}
ઉપરની વ્યાખ્યાને સમકક્ષ
બીજી વ્યાખ્યા આ છે:
વાક્યોનો ગણ $\Gamma$
\emph{મહત્તમ સુસંગત} છે
જો અને માત્ર જો
\begin{enumerate}
\item $\Gamma$ સુસંગત હોય, અને
\item જો $\Gamma \cup \{ !A \}$ સુસંગત હોય,
તો $!A \in \Gamma$ હોય.
\end{enumerate}
બીજા શબ્દોમાં, મહત્તમ સુસંગત
ગણ~$\Gamma$માં પહેલેથી ન
હોય એવું !!{sentence} એ જ $\Gamma$માં ઉમેરીએ,
તો મોટો થયેલો ગણ
અસુસંગત થયા વિના રહેતો નથી.
\end{explain}

\begin{explain}
પૂર્ણતાની સાબિતીમાં મહત્તમ સુસંગત
ગણો મહત્વના છે. કારણ કે
!!{sentence}sનો દરેક સુસંગત
ગણ~$\Gamma$ કોઈ મહત્તમ
સુસંગત ગણ~$\Gamma^*$માં
સમાયેલો છે તેની ખાતરી
આપી શકાય છે; અને
મહત્તમ સુસંગત ગણમાં
દરેક !!{sentence}~$!A$
માટે કાં તો $!A$ અથવા
તેનો નિષેધ $\lnot !A$
હોય છે. ખાસ કરીને
પરમાણુક !!{sentence}s
માટે આ સાચું છે.
આથી !!{constant}s વડે
યોગ્ય રીતે વિસ્તૃત
ભાષામાંના મહત્તમ
સુસંગત ગણમાંથી
એવી !!a{structure}
રચી શકાય છે જેમાં
!!{predicate}sનું અર્થઘટન
$\Gamma^*$માં કયાં
પરમાણુક !!{sentence}s
છે તેના આધારે
વ્યાખ્યાયિત થાય છે.
પછી આ !!{structure}
$\Gamma^*$નાં બધાં
!!{sentence}sને
(અને તેથી $\Gamma$નાં
વાક્યોને પણ) સાચાં
બનાવે છે તેમ બતાવી
શકાય છે. આ છેલ્લી
હકીકતની સાબિતી માટે
$\lnot !A \in
\Gamma^*$ જો અને માત્ર જો
$!A \notin \Gamma^*$,
તેમજ $(!A \lor !B) \in \Gamma^*$
જો અને માત્ર જો
$!A \in \Gamma^*$ અથવા
$!B \in \Gamma^*$,
વગેરે સંબંધો જરૂરી છે.
\end{explain}

\begin{prop}
\ollabel{prop:mcs}
ધારો કે $\Gamma$ મહત્તમ સુસંગત છે. તો:
\begin{enumerate}
\item \ollabel{prop:mcs-prov-in} જો $\Gamma \Proves !A$, તો $!A \in
  \Gamma$.

\item \ollabel{prop:mcs-either-or} કોઈ પણ $!A$ માટે
  કાં તો $!A \in \Gamma$ અથવા $\lnot !A \in \Gamma$.

\tagitem{prvAnd}{\ollabel{prop:mcs-and} $(!A \land !B) \in \Gamma$
  જો અને માત્ર જો $!A \in \Gamma$ અને $!B \in \Gamma$ બંને.}{}

\tagitem{prvOr}{\ollabel{prop:mcs-or} $(!A \lor !B) \in \Gamma$
  જો અને માત્ર જો કાં તો $!A \in \Gamma$ અથવા $!B \in \Gamma$.}{}

\tagitem{prvIf}{\ollabel{prop:mcs-if} $(!A \lif !B) \in \Gamma$
  જો અને માત્ર જો કાં તો $!A \notin \Gamma$ અથવા $!B \in \Gamma$.}{}
\end{enumerate}
\end{prop}

\begin{proof}
આગળના બધા કિસ્સાઓમાં
$\Gamma$ મહત્તમ સુસંગત
છે એમ ધારો.
\begin{enumerate}
\item જો $\Gamma \Proves !A$, તો $!A \in \Gamma$.

ધારો કે $\Gamma \Proves !A$.
પ્રતિપક્ષે $!A \notin \Gamma$
માનો. $\Gamma$ મહત્તમ
સુસંગત હોવાથી
$\Gamma \cup \{!A\}$
અસુસંગત છે; એટલે
$\Gamma \cup \{!A\} \Proves
\lfalse$. હવે
\tagrefs{prfSC/{fol:seq:prv:prop:provability-contr},prfND/{fol:ntd:prv:prop:provability-contr}},
મુજબ $\Gamma$ અસુસંગત
થાય છે. આ $\Gamma$
સુસંગત છે એ
પરિકલ્પનાનો વિરોધાભાસ
છે. તેથી $!A \notin
\Gamma$ હોઈ શકે નહીં;
એટલે $!A \in \Gamma$.

\item કોઈ પણ $!A$ માટે કાં તો $!A \in \Gamma$
અથવા $\lnot !A \in \Gamma$.

પ્રતિપક્ષે માનો કે કોઈ~$!A$
માટે $!A \notin \Gamma$
અને $\lnot !A \notin \Gamma$
બંને છે. $\Gamma$
મહત્તમ સુસંગત હોવાથી
$\Gamma \cup \{!A\}$
અને $\Gamma \cup \{\lnot !A\}$
બંને અસુસંગત છે.
એટલે $\Gamma \cup \{!A\} \Proves \lfalse$
અને $\Gamma \cup
\{\lnot !A\} \Proves \lfalse$.
હવે
\tagrefs{prfSC/{fol:seq:prv:prop:provability-exhaustive},prfND/{fol:ntd:prv:prop:provability-exhaustive}},
મુજબ $\Gamma$ અસુસંગત
થાય છે, જે વિરોધાભાસ છે.
આવું !!a{sentence}~$!A$
હોઈ શકે નહીં; એટલે
દરેક $!A$ માટે
$!A \in \Gamma$ અથવા
$\lnot !A \in \Gamma$.

\tagitem{defAnd}{}{%
\iftag{probAnd}{કસરત.}{%
$(!A \land !B) \in \Gamma$ જો અને માત્ર જો
$!A \in \Gamma$ અને $!B \in \Gamma$ બંને:

આગળની દિશા માટે
$(!A \land !B) \in \Gamma$
ધારો. ત્યારે
$\Gamma \Proves !A \land !B$.
હવે
\tagrefs{prfSC/{fol:seq:prv:prop:provability-land-left},prfND/{fol:ntd:prv:prop:provability-land-left}},
મુજબ $\Gamma \Proves !A$
અને $\Gamma \Proves !B$.
\olref{prop:mcs-prov-in}થી
$!A \in \Gamma$ અને
$!B \in \Gamma$,
જેમ જોઈતું હતું.

ઊલટી દિશા માટે
$!A \in \Gamma$ અને
$!B \in \Gamma$ લો.
ત્યારે $\Gamma \Proves !A$
અને $\Gamma \Proves !B$.
હવે
\tagrefs{prfSC/{fol:seq:prv:prop:provability-land-right},prfND/{fol:ntd:prv:prop:provability-land-right}},
મુજબ $\Gamma \Proves !A \land !B$.
\olref{prop:mcs-prov-in}થી
$(!A \land
!B) \in \Gamma$.}}

\tagitem{defOr}{}{%
\iftag{probOr}{કસરત.}{%
$(!A \lor !B) \in \Gamma$ જો અને માત્ર જો
કાં તો $!A \in \Gamma$ અથવા $!B \in \Gamma$.

આગળની દિશાના પ્રતિપક્ષ
માટે $!A \notin \Gamma$
અને $!B \notin \Gamma$
ધારો. આપણે $(!A \lor
!B) \notin \Gamma$
બતાવવાનું છે. $\Gamma$
મહત્તમ સુસંગત હોવાથી
$\Gamma \cup \{!A\} \Proves \lfalse$
અને $\Gamma \cup \{!B\} \Proves \lfalse$.
હવે
\tagrefs{prfSC/{fol:seq:prv:prop:provability-lor-left},prfND/{fol:ntd:prv:prop:provability-lor-left}},
મુજબ $\Gamma \cup \{(!A \lor !B)\}$
અસુસંગત છે. તેથી
$(!A \lor !B) \notin \Gamma$,
જેમ જોઈતું હતું.

ઊલટી દિશા માટે
$!A \in \Gamma$ અથવા
$!B \in \Gamma$ ધારો.
ત્યારે $\Gamma \Proves !A$
અથવા $\Gamma \Proves !B$.
હવે
\tagrefs{prfSC/{fol:seq:prv:prop:provability-lor-right},prfND/{fol:ntd:prv:prop:provability-lor-right}},
મુજબ $\Gamma \Proves !A \lor !B$.
\olref{prop:mcs-prov-in}થી
$(!A \lor !B) \in \Gamma$,
જેમ જોઈતું હતું.}}

\tagitem{defIf}{}{%
\iftag{probIf}{કસરત.}{%
$(!A \lif !B) \in \Gamma$ જો અને માત્ર જો
કાં તો $!A \notin \Gamma$ અથવા $!B \in \Gamma$:

આગળની દિશા માટે
$(!A \lif !B) \in \Gamma$
લો અને પ્રતિપક્ષે
$!A \in \Gamma$ તથા
$!B \notin \Gamma$ ધારો.
આ પરિકલ્પનાઓથી
$\Gamma \Proves !A \lif !B$
અને $\Gamma \Proves !A$.
હવે
\tagrefs{prfSC/{fol:seq:prv:prop:provability-mp},prfND/{fol:ntd:prv:prop:provability-mp}},
મુજબ $\Gamma \Proves !B$.
પરંતુ પછી
\olref{prop:mcs-prov-in}થી
$!B \in \Gamma$,
જે $!B \notin \Gamma$
એ ધારણાનો વિરોધાભાસ છે.

ઊલટી દિશા માટે
પહેલાં $!A \notin \Gamma$
હોય તે કિસ્સો લો.
\olref{prop:mcs-either-or}થી
$\lnot !A \in \Gamma$
અને તેથી
$\Gamma \Proves \lnot !A$.
હવે
\tagrefs{prfSC/{fol:seq:prv:prop:provability-lif},prfND/{fol:ntd:prv:prop:provability-lif}},
મુજબ $\Gamma \Proves !A \lif !B$.
ફરી \olref{prop:mcs-prov-in}થી
$(!A \lif !B) \in \Gamma$
મળે, જેમ જોઈતું હતું.

હવે $!B \in \Gamma$
હોય તે કિસ્સો લો.
ત્યારે $\Gamma \Proves !B$,
અને
\tagrefs{prfSC/{fol:seq:prv:prop:provability-lif},prfND/{fol:ntd:prv:prop:provability-lif}},
મુજબ $\Gamma \Proves !A \lif !B$.
\olref{prop:mcs-prov-in}થી
$(!A \lif !B) \in \Gamma$.}}
\end{enumerate}
\end{proof}

\begin{prob}
\olref[fol][com][mcs]{prop:mcs}ની સાબિતી પૂરી કરો.
\end{prob}

\end{document}
